% Two editorial finite constructions; neither proves a historical truth criterion.
\begin{annexitem}{annex:p15.wm.causal-ambiguity}{thm:p15-wm-causal-ambiguity}{Agreement and covariation can leave the cause unsettled}
\begin{proof}
The two contexts have intersection \(\{A\}\), since the five circumstance
names are distinct. The first contains \(C\) and the second contains \(E\),
so the stipulated occurrence condition gives \(a\) in both.
After removing \(A\), the remaining contexts are \(\{B,C\}\) and
\(\{D,E\}\). They still contain \(C\) and \(E\), respectively, so
\(a\) still occurs. In fact, inserting or removing \(A\) in any context
cannot affect an occurrence condition referring only to \(C\) and \(E\).
This verifies the counterexample using the complete declared context,
without postulating an omitted shared circumstance.

In the direct model, substitution of \(X=U\) into \(Y=X\) gives \(Y=U\).
In the common-source model \(X=U\) and \(Y=U\) already.
Consequently both give \((X,Y)=(0,0)\) when \(U=0\) and
\((X,Y)=(1,1)\) when \(U=1\), each with probability \(1/2\);
the other two pairs have probability zero.

After the declared replacement, the direct model has \(X=0\) and
the unchanged equation \(Y=X\), hence \(Y=0\) for either value of \(U\).
The common-source model has \(X=0\) but retains \(Y=U\), hence
probability \(1/2\) for each value of \(Y\).
The resulting \(Y\)-laws therefore differ although the unreplaced
observational laws agree. These computations use exactly the specified
replacement and unchanged exogenous law; no claim that observational
equality fixes a physical intervention mechanism is needed.
\end{proof}
\end{annexitem}

\begin{annexitem}{annex:p15.wm.finite-continuation}{thm:p15-wm-finite-continuation}{Exact finite fit does not determine extrapolation}
\begin{proof}
All denominators \(x_i-x_j\) appearing in the sum are nonzero by the
distinctness assumption, so each term is a well-defined real polynomial.
At \(x=x_k\), the term indexed by \(i=k\) has every factor equal to one.
Every other term has the factor \(x_k-x_k=0\), because its product includes
\(j=k\). Thus \(p(x_k)=y_k\). For \(m=1\) the single empty product is one,
and the same evaluation gives the constant polynomial \(p(x)=y_1\).

The product \(\prod_i(x-x_i)\) vanishes at every \(x_k\), so adding
any real multiple of it does not change those fitted values.
If \(t\) is not an observed abscissa, set
\[
 R=\prod_{i=1}^m(t-x_i).
\]
Every factor is nonzero, hence \(R\ne0\).
For a prescribed real value \(z\), choose \(c=(z-p(t))/R\).
Then \(p_c(t)=p(t)+cR=z\). This proves the claim for every such \(t\),
not merely for one extrapolation point.

For the displayed two-point example, \(p_0(0)=p_1(0)=0\) and
\(p_0(1)=p_1(1)=1\), while \(p_0(2)=2\) and
\(p_1(2)=2+2(2-1)=4\), verifying the stated numerical comparison.
No probabilistic noise, degree bound or sampling premise is assumed.
\end{proof}
\end{annexitem}
